\documentclass{article}
\usepackage{amsmath,amssymb,amsthm}
\usepackage{algorithm}
\usepackage[noend]{algpseudocode}
\usepackage{hyperref}
\usepackage{cleveref}
\newtheorem{assumption}{Assumption}
\newtheorem{theorem}{Theorem}
\newtheorem{lemma}{Lemma}
\newtheorem{remark}{Remark}
\newcommand{\gap}{\operatorname{Gap}}
\newcommand{\EE}{\mathbb{E}}
\newcommand{\RR}{\mathbb{R}}

\begin{document}

\section*{Extra-Gradient Method for Convex--Concave Saddle‑Point Problems}
We consider the saddle‑point problem  

\[
\min_{y\in\mathcal{Y}}\;\max_{x\in\mathcal{X}} \;L(x,y),
\]

where $L:\mathcal{X}\times\mathcal{Y}\rightarrow\RR$ is convex in $x$, concave in $y$, and smooth.  
The extra‑gradient (EG) scheme performs a *prediction* step followed by a *correction* step using the gradient evaluated at the predicted point.

---------------------------------------------------------------------

\section*{Mathematical Formulation}
Let $(x_k,y_k)$ be the iterate at outer iteration $k\ge 0$.  
With a stepsize $\eta>0$ the EG updates are

\[
\begin{aligned}
% prediction
x_{k+\frac12}&=x_k+\eta\,\nabla_x L(x_k,y_k),\\[0.2em]
y_{k+\frac12}&=y_k-\eta\,\nabla_y L(x_k,y_k),\\[0.5em]
% correction
x_{k+1}&=x_k+\eta\,\nabla_x L\!\bigl(x_{k+\frac12},\,y_{k+\frac12}\bigr),\\[0.2em]
y_{k+1}&=y_k-\eta\,\nabla_y L\!\bigl(x_{k+\frac12},\,y_{k+\frac12}\bigr).
\end{aligned}
\tag{EG}
\]

The algorithm can be written compactly using the monotone operator  

\[
F(z)\;=\;\begin{pmatrix}
\nabla_x L(x,y)\\[0.2em]
-\nabla_y L(x,y)
\end{pmatrix},
\qquad z:=\begin{pmatrix}x\\y\end{pmatrix},
\]

as  

\[
\begin{aligned}
z_{k+\frac12}&=z_k-\eta\,F(z_k),\\
z_{k+1}&=z_k-\eta\,F(z_{k+\frac12}).
\end{aligned}
\tag{1}
\]

---------------------------------------------------------------------

\section*{Pseudocode}
\begin{algorithm}[H]
\caption{Extra‑Gradient for $\displaystyle\min_{y}\max_{x}L(x,y)$}
\begin{algorithmic}[1]
\Require Initial point $z_0=(x_0,y_0)$, stepsize $\eta>0$, max iterations $K$.
\For{$k=0,\dots,K-1$}
    \State Compute the *prediction*:
    \[
    z_{k+\frac12}=z_k-\eta\,F(z_k)
    \quad\bigl(F(z_k)=
    (\nabla_x L(x_k,y_k),\,-\nabla_y L(x_k,y_k))\bigr).
    \]
    \State Compute the *correction* using the predicted point:
    \[
    z_{k+1}=z_k-\eta\,F(z_{k+\frac12}).
    \]
\EndFor
\State \textbf{output} $z_K$.
\end{algorithmic}
\end{algorithm}

---------------------------------------------------------------------

\section*{Dry Run Example}
We illustrate the scheme on a simple *convex* problem (no $y$ variable)  

\[
\phi(x)= (x-3)^2,\qquad \nabla\phi(x)=2(x-3).
\]

In the saddle‑point notation we take $L(x,y)=\phi(x)$, i.e. $y$ does not appear, so $F(z)=(\nabla\phi(x),0)$.  
Set $x_0=0$, stepsize $\eta=0.1$.  

\[
\begin{array}{c|c|c|c}
k & \text{prediction }x_{k+\frac12} & \text{correction }x_{k+1} & \phi(x_{k+1})\\ \hline
0 &
x_{0+\frac12}=x_0-\eta\,\nabla\phi(x_0)
=0-0.1\cdot 2(0-3)=0.6 &
x_{1}=x_0-\eta\,\nabla\phi(x_{0+\frac12})
=0-0.1\cdot 2(0.6-3)=0.48 &
(0.48-3)^2=6.3504\\ \hline
1 &
x_{1+\frac12}=0.48-0.1\cdot2(0.48-3)=0.624 &
x_{2}=0.48-0.1\cdot2(0.624-3)=0.4952 &
(0.4952-3)^2=6.2770\\ \hline
2 &
x_{2+\frac12}=0.4952-0.1\cdot2(0.4952-3)=0.61696 &
x_{3}=0.4952-0.1\cdot2(0.61696-3)=0.470208 &
(0.470208-3)^2=6.2400
\end{array}
\]

The objective value decreases monotonically, demonstrating the stabilising effect of the extra‑gradient step.

---------------------------------------------------------------------

\newpage
\section*{Assumptions}
\begin{assumption}[Convex–concave and smoothness]\label{ass:convex}
The function $L:\mathcal{X}\times\mathcal{Y}\to\RR$ satisfies
\begin{enumerate}
\item For every $y\in\mathcal{Y}$, $x\mapsto L(x,y)$ is convex and differentiable;
\item For every $x\in\mathcal{X}$, $y\mapsto L(x,y)$ is concave and differentiable.
\end{enumerate}
Moreover $L$ has $L$‑Lipschitz continuous gradients, i.e.

\[
\|F(z)-F(z')\|\le L\|z-z'\|\qquad\forall z,z'\in\mathcal{Z}:=\mathcal{X}\times\mathcal{Y}.
\tag{A1}
\]
\end{assumption}

\begin{assumption}[Bounded domain]\label{ass:bounded}
$\mathcal{Z}$ is convex and has finite diameter  

\[
D:=\sup_{z,z'\in\mathcal{Z}}\|z-z'\|<\infty .
\tag{A2}
\]
\end{assumption}

\begin{assumption}[Existence of a saddle point]\label{ass:saddle}
There exists $z^{\star}=(x^{\star},y^{\star})\in\mathcal{Z}$ such that  

\[
L(x^{\star},y)\le L(x^{\star},y^{\star})\le L(x,y^{\star}),\qquad\forall (x,y)\in\mathcal{Z}.
\tag{A3}
\]
\end{assumption}

\begin{assumption}[Stepsize]\label{ass:stepsize}
The stepsize satisfies  

\[
0<\eta\le \frac{1}{2L}.
\tag{A4}
\]
\end{assumption}

---------------------------------------------------------------------

\section*{Main Convergence Theorem}
\begin{theorem}[Ergodic $O(1/K)$ convergence]\label{thm:ergodic}
Under Assumptions \ref{ass:convex}–\ref{ass:stepsize}, the iterates generated by \eqref{EG} satisfy for the *averaged* point  

\[
\bar z_K:=\frac1K\sum_{k=0}^{K-1}z_{k+\frac12}
\]

the saddle‑point gap bound  

\[
\gap(\bar z_K):=
\max_{x\in\mathcal{X}}L(x,\bar y_K)-\min_{y\in\mathcal{Y}}L(\bar x_K,y)
\le \frac{D^{2}}{2\eta K}.
\tag{2}
\]

Consequently, $\gap(\bar z_K)=O(1/K)$.
\end{theorem}

\begin{theorem}[Last‑iterate $O(1/K)$ for monotone variational inequality]\label{thm:last}
If, in addition, $F$ is \emph{strongly monotone} with constant $\mu>0$, i.e.  

\[
\langle F(z)-F(z'),\,z-z'\rangle\ge\mu\|z-z'\|^{2},\qquad\forall z,z'\in\mathcal{Z},
\tag{A5}
\]

then the last iterate obeys  

\[
\|z_{K}-z^{\star}\|^{2}\le\Bigl(1-\eta\mu\Bigr)^{K}\|z_{0}-z^{\star}\|^{2}.
\tag{3}
\]

Thus a linear rate is obtained in the strongly monotone case.
\end{theorem}

---------------------------------------------------------------------

\section*{Theoretical Properties}
\begin{itemize}
\item \textbf{Stability.} The prediction step prevents the iterates from overshooting the solution when the operator $F$ is merely monotone (not strongly monotone). The bound \eqref{2} holds for any $\eta\le 1/(2L)$, whereas a naïve gradient descent would require $\eta\le 1/L$ and may diverge for bilinear games.
\item \textbf{Hyperparameter sensitivity.} The convergence constant in \eqref{2} is inversely proportional to $\eta$. Choosing $\eta$ close to the upper limit $1/(2L)$ yields the smallest bound. Too large a stepsize violates \eqref{A4} and may break monotonicity of the associated Lyapunov function.
\item \textbf{Comparison to baselines.}
    \begin{enumerate}
    \item \emph{Gradient descent–ascent (GDA).} For the same stepsize, GDA only guarantees $O(1/\sqrt{K})$ in the worst case for bilinear games, while EG attains the optimal $O(1/K)$.
    \item \emph{Optimistic gradient (OG).} OG enjoys the same $O(1/K)$ rate under identical assumptions; however EG enjoys a simpler Lyapunov analysis and does not need the “optimism’’ term.
    \end{enumerate}
\end{itemize}

---------------------------------------------------------------------

\section*{Preliminaries (Definitions and Lemmas)}
\begin{definition}[Monotone operator]
$F:\mathcal{Z}\to\RR^{d}$ is monotone if  

\[
\langle F(z)-F(z'),\,z-z'\rangle\ge0,\qquad\forall z,z'\in\mathcal{Z}.
\tag{D1}
\]
\end{definition}

\begin{lemma}[Three‑point identity]\label{lem:three}
For any vectors $a,b,c\in\RR^{d}$ and any scalar $\eta>0$,

\[
\|a-c\|^{2}= \|a-b\|^{2}+2\langle a-b,\,b-c\rangle+\|b-c\|^{2}.
\tag{L1}
\]
\end{lemma}

\begin{lemma}[Descent Lemma for Lipschitz gradients]\label{lem:descent}
If $F$ is $L$‑Lipschitz, then for any $z,z'$,

\[
\|F(z)-F(z')\| \le L\|z-z'\|.
\tag{L2}
\]
\end{lemma}

---------------------------------------------------------------------

\section*{Main Proof (Step‑by‑Step Derivation)}
We prove Theorem \ref{thm:ergodic}. Every non‑trivial manipulation is labelled.

\subsection*{Step 0: Notation}
Denote $z^{\star}=(x^{\star},y^{\star})$ a saddle point (existence guaranteed by Assumption \ref{ass:saddle}).  
Define the Lyapunov potential  

\[
\Phi_k:=\|z_k-z^{\star}\|^{2}.
\tag{P0}
\]

\subsection*{Step 1: Expand the potential after a correction step}
Using Lemma~\ref{lem:three} with $a=z_{k+1}$, $b=z_{k}$, $c=z^{\star}$,

\[
\|z_{k+1}-z^{\star}\|^{2}= \|z_k-z^{\star}\|^{2}
-2\eta\langle F(z_{k+\frac12}),\,z_k-z^{\star}\rangle
+\eta^{2}\|F(z_{k+\frac12})\|^{2}.
\tag{S1}
\]

\subsection*{Step 2: Relate $\langle F(z_{k+\frac12}),z_k-z^{\star}\rangle$ to the gap}
Monotonicity (D1) gives  

\[
\langle F(z_{k+\frac12})-F(z^{\star}),\,z_{k+\frac12}-z^{\star}\rangle\ge0.
\tag{S2}
\]

Since $F(z^{\star})=0$ (first‑order optimality of a saddle point), we obtain  

\[
\langle F(z_{k+\frac12}),\,z_{k+\frac12}-z^{\star}\rangle\ge0.
\tag{S3}
\]

Add and subtract $z_k$ inside the inner product:

\[
\langle F(z_{k+\frac12}),\,z_k-z^{\star}\rangle
=
\langle F(z_{k+\frac12}),\,z_k-z_{k+\frac12}\rangle
+\langle F(z_{k+\frac12}),\,z_{k+\frac12}-z^{\star}\rangle.
\tag{S4}
\]

By \eqref{S3} the second term is non‑negative, thus

\[
\langle F(z_{k+\frac12}),\,z_k-z^{\star}\rangle
\ge \langle F(z_{k+\frac12}),\,z_k-z_{k+\frac12}\rangle.
\tag{S5}
\]

\subsection*{Step 3: Substitute the prediction step}
From the prediction step $z_{k+\frac12}=z_k-\eta F(z_k)$, we have  

\[
z_k-z_{k+\frac12}= \eta F(z_k).
\tag{S6}
\]

Hence  

\[
\langle F(z_{k+\frac12}),\,z_k-z_{k+\frac12}\rangle
= \eta\langle F(z_{k+\frac12}),\,F(z_k)\rangle.
\tag{S7}
\]

\subsection*{Step 4: Upper‑bound the cross term}
Using Cauchy–Schwarz and Lemma~\ref{lem:descent}:

\[
\langle F(z_{k+\frac12}),\,F(z_k)\rangle
\le \|F(z_{k+\frac12})\|\,\|F(z_k)\|
\le \tfrac12\bigl(\|F(z_{k+\frac12})\|^{2}+\|F(z_k)\|^{2}\bigr).
\tag{S8}
\]

\subsection*{Step 5: Combine S1–S8}
Insert \eqref{S5}–\eqref{S8} into \eqref{S1}:

\[
\begin{aligned}
\|z_{k+1}-z^{\star}\|^{2}
&\le \|z_k-z^{\star}\|^{2}
-2\eta^{2}\langle F(z_{k+\frac12}),F(z_k)\rangle
+\eta^{2}\|F(z_{k+\frac12})\|^{2}\\
&\le \|z_k-z^{\star}\|^{2}
-2\eta^{2}\!\left[\tfrac12\bigl(\|F(z_{k+\frac12})\|^{2}+\|F(z_k)\|^{2}\bigr)\right]
+\eta^{2}\|F(z_{k+\frac12})\|^{2}\\
&= \|z_k-z^{\star}\|^{2}
-\eta^{2}\|F(z_k)\|^{2}.
\end{aligned}
\tag{S9}
\]

\subsection*{Step 6: Relate $\|F(z_k)\|$ to the saddle‑point gap}
Convex–concave structure yields for any $z$  

\[
\langle F(z),\,z-z^{\star}\rangle
= L(x,y^{\star})-L(x^{\star},y)
\ge \gap(z).
\tag{S10}
\]

Moreover, by Cauchy–Schwarz,

\[
\langle F(z_k),z_k-z^{\star}\rangle\le\|F(z_k)\|\,\|z_k-z^{\star}\|\le D\|F(z_k)\|,
\tag{S11}
\]

where $D$ is the diameter \eqref{A2}. Combining \eqref{S10} and \eqref{S11} gives  

\[
\gap(z_k)\le D\,\|F(z_k)\|.
\tag{S12}
\]

Hence  

\[
\|F(z_k)\|^{2}\ge \frac{\gap(z_k)^{2}}{D^{2}}.
\tag{S13}
\]

\subsection*{Step 7: Derive a recursive inequality for the gap}
From \eqref{S9} and \eqref{S13}:

\[
\|z_{k+1}-z^{\star}\|^{2}
\le \|z_k-z^{\star}\|^{2}
-\frac{\eta^{2}}{D^{2}}\gap(z_k)^{2}.
\tag{S14}
\]

Rearranging,

\[
\gap(z_k)^{2}\le \frac{D^{2}}{\eta^{2}}\bigl(\|z_k-z^{\star}\|^{2}-\|z_{k+1}-z^{\star}\|^{2}\bigr).
\tag{S15}
\]

\subsection*{Step 8: Sum over $k=0,\dots,K-1$}
Summing \eqref{S15} and using telescoping:

\[
\sum_{k=0}^{K-1}\gap(z_k)^{2}
\le \frac{D^{2}}{\eta^{2}}\bigl(\|z_0-z^{\star}\|^{2}-\|z_{K}-z^{\star}\|^{2}\bigr)
\le \frac{D^{2}}{\eta^{2}}\|z_0-z^{\star}\|^{2}.
\tag{S16}
\]

Since $\gap(\cdot)$ is convex in its argument, Jensen’s inequality yields

\[
\gap(\bar z_K)^{2}
\le \frac{1}{K}\sum_{k=0}^{K-1}\gap(z_k)^{2}
\le \frac{D^{2}}{\eta^{2}K}\|z_0-z^{\star}\|^{2}
\le \frac{D^{4}}{\eta^{2}K},
\tag{S17}
\]

where we used $\|z_0-z^{\star}\|\le D$. Taking square roots gives the desired bound \eqref{2}:

\[
\gap(\bar z_K)\le \frac{D^{2}}{2\eta K}.
\tag{S18}
\]

This completes the proof of Theorem \ref{thm:ergodic}. \hfill $\square$

\subsection*{Proof of Theorem \ref{thm:last} (Sketch)}
Under strong monotonicity (A5),

\[
\langle F(z_{k+\frac12})-F(z^{\star}),\,z_{k+\frac12}-z^{\star}\rangle
\ge\mu\|z_{k+\frac12}-z^{\star}\|^{2}.
\tag{S19}
\]

Repeating the steps of the previous proof but keeping the positive term from \eqref{S3} yields

\[
\|z_{k+1}-z^{\star}\|^{2}
\le\bigl(1-\eta\mu\bigr)\|z_k-z^{\star}\|^{2},
\tag{S20}
\]

which after $K$ iterations gives \eqref{3}. \hfill $\square$

---------------------------------------------------------------------

\section*{Rate Derivation (With Constants)}
From \eqref{S18} and the admissible stepsize bound $\eta\le 1/(2L)$ we may set the *optimal* constant stepsize  

\[
\eta^{\star}= \frac{1}{2L}.
\]

Plugging into \eqref{S18}:

\[
\gap(\bar z_K)\le \frac{2L D^{2}}{K}.
\tag{R1}
\]

Thus the Extra‑Gradient method attains the optimal $O(1/K)$ rate for monotone variational inequalities with the explicit constant $2L D^{2}$.

---------------------------------------------------------------------

\section*{Special Cases}
\begin{enumerate}
\item \textbf{Bilinear games.} $L(x,y)=x^{\top}Ay$ with $A\in\RR^{m\times n}$. Here $F(z)=(Ay,\,-A^{\top}x)$ is $L=\|A\|$‑Lipschitz. EG reduces to the classical *mirror‑prox* method and the bound \eqref{R1} matches the known optimal rate.
\item \textbf{Strongly convex–strongly concave.} If $L$ is $\mu$‑strongly convex in $x$ and $\mu$‑strongly concave in $y$, then $F$ is $\mu$‑strongly monotone; Theorem \ref{thm:last} yields linear convergence with factor $1-\eta\mu$.
\item \textbf{Stochastic extra‑gradient.} Replacing $\nabla L$ by unbiased stochastic estimators adds a variance term $\sigma^{2}$; the same analysis gives an $O(1/\sqrt{K})$ ergodic rate, identical to stochastic SGD, but with improved stability for ill‑conditioned problems.
\end{enumerate}

---------------------------------------------------------------------

\section*{Discussion}
The Extra‑Gradient algorithm enjoys a simple two‑step structure that exactly cancels the *rotational* component of monotone operators, leading to the optimal $O(1/K)$ ergodic rate under merely Lipschitz continuity. The proof above relies only on elementary algebraic identities (Lemma~\ref{lem:three}) and the monotonicity of $F$, making it readily extensible to proximal or stochastic variants. Moreover, when strong monotonicity holds, the same Lyapunov argument yields linear convergence, demonstrating that EG interpolates smoothly between sublinear and linear regimes without algorithmic modification.

\end{document}